\documentclass[a4paper,11pt]{article}

\newcommand{\TPnumero}{Lab 6 -- ARCH/GARCH}
% =========================================================
% Tutorial / lab preamble — Time Series Analysis module
% article class + fancyhdr + tcolorbox + listings (English)
% Author : Aymen Ben Brik (aymen.benbrik@esprit.tn)
% =========================================================

\usepackage[utf8]{inputenc}
\usepackage[T1]{fontenc}
\usepackage[english]{babel}
\usepackage{lmodern}
\usepackage{amsmath, amssymb, amsthm, mathtools, bm}
\usepackage{graphicx}
\usepackage{booktabs}
\usepackage{array}
\usepackage{multirow}
\usepackage{colortbl}
\usepackage{xcolor}
\usepackage{tikz}
\usetikzlibrary{positioning, arrows.meta, calc, shapes, fit, backgrounds, matrix}
\usepackage{listings}
\usepackage[most]{tcolorbox}
\usepackage{geometry}
\geometry{a4paper, margin=2cm, headheight=15pt}
\usepackage{fancyhdr}
\usepackage[hidelinks]{hyperref}
\usepackage{enumitem}

% --- Colours ---
\definecolor{espritBlue}{RGB}{30, 60, 110}
\definecolor{espritAccent}{RGB}{200, 80, 30}
\definecolor{boxEx}{RGB}{30, 60, 110}
\definecolor{boxQ}{RGB}{180, 120, 0}
\definecolor{boxC}{RGB}{30, 100, 60}
\definecolor{boxAtt}{RGB}{200, 80, 30}

% --- Listings ---
\definecolor{codebg}{RGB}{248,248,248}
\definecolor{codeframe}{RGB}{220,220,220}
\definecolor{keyword}{RGB}{0,82,155}
\definecolor{comment}{RGB}{88,110,117}
\definecolor{string}{RGB}{133,153,0}

\lstdefinestyle{pythonstyle}{
  language=Python,
  basicstyle=\ttfamily\small,
  keywordstyle=\color{keyword}\bfseries,
  commentstyle=\color{comment}\itshape,
  stringstyle=\color{string},
  backgroundcolor=\color{codebg},
  frame=single,
  rulecolor=\color{codeframe},
  frameround=tttt,
  breaklines=true,
  showstringspaces=false,
  tabsize=2
}
\lstdefinestyle{Rstyle}{
  language=R,
  basicstyle=\ttfamily\small,
  keywordstyle=\color{keyword}\bfseries,
  commentstyle=\color{comment}\itshape,
  stringstyle=\color{string},
  backgroundcolor=\color{codebg},
  frame=single,
  rulecolor=\color{codeframe},
  frameround=tttt,
  breaklines=true,
  showstringspaces=false,
  tabsize=2
}
\lstset{style=pythonstyle}

% --- Common macros (configurable path) ---
\providecommand{\macrosPath}{../../preamble/macros.tex}
\input{\macrosPath}

% --- Exercise counter ---
\newcounter{excounter}
\renewcommand{\theexcounter}{\arabic{excounter}}

\newtcolorbox[use counter=excounter]{exercise}[1][]{
  enhanced, breakable,
  colback=boxEx!5, colframe=boxEx, fonttitle=\bfseries,
  title={Exercise~\theexcounter\ifstrempty{#1}{}{ -- #1}},
  arc=2pt, boxrule=0.6pt
}
\newtcolorbox{question}[1][]{
  enhanced, breakable,
  colback=boxQ!5, colframe=boxQ, fonttitle=\bfseries,
  title={Question\ifstrempty{#1}{}{ -- #1}},
  arc=2pt, boxrule=0.5pt
}
\newtcolorbox{solution}[1][]{
  enhanced, breakable,
  colback=boxC!5, colframe=boxC, fonttitle=\bfseries,
  title={Solution\ifstrempty{#1}{}{ -- #1}},
  arc=2pt, boxrule=0.5pt
}
\newtcolorbox{warning}[1][]{
  enhanced, breakable,
  colback=boxAtt!5, colframe=boxAtt, fonttitle=\bfseries,
  title={Warning\ifstrempty{#1}{}{ -- #1}},
  arc=2pt, boxrule=0.5pt
}

% --- Headers / footers ---
\pagestyle{fancy}
\fancyhf{}
\fancyhead[L]{\textsc{\moduleNom} -- \auteurNom}
\fancyhead[R]{\TPnumero}
\fancyfoot[C]{\thepage}
\fancyfoot[L]{\auteurInstitution}
\fancyfoot[R]{\auteurEmail}
\renewcommand{\headrulewidth}{0.4pt}
\renewcommand{\footrulewidth}{0.2pt}

% --- Placeholder ---
\providecommand{\TPnumero}{Lab}


\title{Lab 6: ARCH effects, GARCH fitting and Value-at-Risk\\
\large Volatility clustering, Engle LM, GARCH$(1,1)$ and forecasting}
\author{\auteurNom\\\auteurEmail\\\auteurInstitution}
\date{\anneeUniversitaire}

\begin{document}
\maketitle

\section*{Goals}
\begin{itemize}
  \item Detect volatility clustering visually and via the Engle LM
        test.
  \item Estimate ARCH$(q)$ and GARCH$(p, q)$ models by maximum
        likelihood.
  \item Validate the fit through standardised-residual diagnostics.
  \item Produce conditional volatility forecasts and a one-day
        Value-at-Risk.
\end{itemize}

\section*{Setup}
\begin{itemize}
  \item Python 3 with \texttt{numpy}, \texttt{pandas},
        \texttt{matplotlib}, \texttt{statsmodels} and
        \texttt{arch}.
  \item \texttt{numpy.random.seed(42)}.
  \item Dataset: \texttt{data/synthetic\_returns.csv} -- $T = 1500$
        synthetic daily returns generated from a GARCH$(1, 1)$ with
        $(\alpha_0, \alpha_1, \beta_1) = (0.05, 0.10, 0.85)$ and
        Gaussian innovations.
\end{itemize}

\begin{warning}
Install the \texttt{arch} package (Kevin Sheppard) once via
\texttt{pip install arch}. The main entry point is
\texttt{from arch import arch\_model}.
\end{warning}

% =========================================================
\begin{exercise}[Visualising volatility clustering]
\textbf{(1)} Load \texttt{data/synthetic\_returns.csv}. Plot the
return series $r_t$ versus $t$. Comment: do you observe quiet
periods alternating with turbulent ones?

\textbf{(2)} Compute and plot the sample ACF of $r_t$ up to
lag~$20$. The returns should look (almost) uncorrelated.

\textbf{(3)} Plot the sample ACF of $r_t^2$ up to lag~$20$. Compare
with the previous plot and comment in 3 lines.
\end{exercise}

% =========================================================
\begin{exercise}[Engle's LM test]
\textbf{(1)} Use
\texttt{statsmodels.stats.diagnostic.het\_arch} on $r_t$ at lags
$5, 10, 20$. Build a small table with columns: lag, LM statistic,
$\chi^2_{q, 0.95}$ critical value, $p$-value, decision.

\textbf{(2)} Conclude: is there evidence of ARCH effects in the
data? Why is this consistent with what the ACF of $r_t^2$ already
showed?
\end{exercise}

% =========================================================
\begin{exercise}[ARCH$(q)$ fitting]
Use \texttt{arch\_model(returns, vol="ARCH", p=q,
mean="Constant")}. Note: in the \texttt{arch} package the
``\texttt{p}'' parameter is the ARCH order.

\textbf{(1)} Fit ARCH$(1)$, ARCH$(3)$ and ARCH$(5)$. Report
$\widehat\alpha_0$, the $\widehat\alpha_i$, the log-likelihood and
the AIC for each.

\textbf{(2)} Compute $\sum_i \widehat\alpha_i$ for each. Is it less
than~$1$ (covariance-stationarity)?

\textbf{(3)} Why does ARCH$(5)$ have only a marginal AIC gain over
ARCH$(3)$? Hint: the true model has volatility memory not directly
captured by simple lagged squared shocks.
\end{exercise}

% =========================================================
\begin{exercise}[GARCH$(1, 1)$ fitting]
\textbf{(1)} Fit GARCH$(1, 1)$ with constant mean and Gaussian
innovations:
\[
  r_t = \mu + \varepsilon_t,
  \qquad
  \sigma_t^2 = \alpha_0
              + \alpha_1\, \varepsilon_{t-1}^2
              + \beta_1\, \sigma_{t-1}^2.
\]
Report $\widehat\mu$, $\widehat\alpha_0$, $\widehat\alpha_1$,
$\widehat\beta_1$, the log-likelihood and the AIC.

\textbf{(2)} Compute $\widehat\alpha_1 + \widehat\beta_1$ and the
implied unconditional variance
$\widehat\sigma^2 = \widehat\alpha_0 / (1 - \widehat\alpha_1 -
\widehat\beta_1)$. Compare with the empirical sample variance.

\textbf{(3)} Compare AIC: GARCH$(1, 1)$ vs.\ ARCH$(5)$. Comment:
is the parsimonious GARCH$(1, 1)$ preferred?
\end{exercise}

% =========================================================
\begin{exercise}[Standardised-residual diagnostics]
Compute $\widehat z_t = \widehat\varepsilon_t / \widehat\sigma_t$
from the GARCH$(1, 1)$ fit.

\textbf{(1)} Plot the ACF of $\widehat z_t$ and of
$\widehat z_t^2$ up to lag~$20$.

\textbf{(2)} Run Engle's LM test on $\widehat z_t^2$ at the same
lags as Exercise~$2$. The model is well-specified if it does
\emph{not} reject anymore.

\textbf{(3)} Plot a QQ-plot of $\widehat z_t$ vs.\ $\mathcal N(0, 1)$.
Comment on the marginal fit.
\end{exercise}

% =========================================================
\begin{exercise}[Volatility forecast and VaR]
\textbf{(1)} On the GARCH$(1, 1)$ fit, compute the one-step
forecast $\widehat\sigma_{T+1}$ and the iterated $h$-step forecasts
$\widehat\sigma_{T+h}^2$ for $h = 1, \dots, 30$. Plot
$\widehat\sigma_{T+h}^2$ together with the unconditional variance
$\widehat\sigma^2$.

\textbf{(2)} Verify numerically the mean-reversion formula
\[
  \widehat\sigma_{T+h}^2
  = \widehat\sigma^2
  + (\widehat\alpha_1 + \widehat\beta_1)^{h-1}
        \bigl(\widehat\sigma_{T+1}^2 - \widehat\sigma^2\bigr).
\]

\textbf{(3)} Compute the one-day Value-at-Risk at level $\alpha
\in \{1\%, 5\%\}$:
\[
  \mathrm{VaR}_{T+1}(\alpha)
  = -\bigl(\widehat\mu + z_\alpha\, \widehat\sigma_{T+1}\bigr).
\]
Compare with the unconditional VaR using the sample~$\sigma$.
Which one is tighter today?
\end{exercise}

\bigskip
\begin{warning}
$z_\alpha = \Phi^{-1}(\alpha)$ is computed via
\texttt{scipy.stats.norm.ppf(0.01)} $= -2.326$ and
\texttt{norm.ppf(0.05)} $= -1.645$.
\end{warning}

\end{document}
